\subsection{Código Fuente}

% --- Intro ---
\begin{frame}{Código Fuente}
    El código se resalta automáticamente usando \texttt{minted}.
\end{frame}

% --- Buenas Prácticas ---
\begin{frame}{Buenas Prácticas}
    \begin{TipP}
        \textbf{Recomendación:} Se recomienda mantener el código en archivos
        \texttt{.cpp} separados dentro de la carpeta \texttt{code/} y usar
        \texttt{\textbackslash twocolumnminted} para incluirlos. Esto asegura que
        el código sea compilable y testeable independientemente del documento.
    \end{TipP}
\end{frame}



% --- Ejemplo con archivo externo ---
%% NOTA: [fragile] es necesario para meter código (minted/verbatim)
\begin{frame}[fragile]{Ejemplo: Archivo externo}
    Ejemplo incluyendo el archivo \texttt{code/template.cpp} directamente:

    % Inserta el archivo code/template.cpp 
    \twocolumnminted{code/template.cpp}{cpp}{20}
\end{frame}

% --- Ejemplo con snippet inline ---
%% NOTA: [fragile] es necesario para meter código (minted/verbatim)
\begin{frame}[fragile]{Ejemplo: Código inline}
    También es posible escribir código directamente en el archivo \texttt{.tex} usando el entorno \texttt{minted}, útil para snippets cortos:

    \begin{minted}{cpp}
    // Ejemplo de código inline
    void solve() {
        int n; cin >> n;
        cout << n * 2 << endl;
    }
    \end{minted}
\end{frame}
